% from ega3-1-fr.tex
\oldpage[III]{96}
Si $X$ est un préschéma dont l'espace sous-jacent est noethérien, ou un
schéma quasi-compact, il résulte de (I, 9.3.1) qu'en posant comme d'ordinaire
$U_{i_0i_1\ldots i_p}=\bigcap_{k=0}^pU_{i_k}$, on a, à un isomorphisme
canonique près,
\[
\Gamma(U_{i_0i_1\ldots i_p},\mathcal F(*))
=H^0(U_{i_0i_1\ldots i_p},\mathcal F(*))
=M_{f_{i_0}f_{i_1}\cdots f_{i_p}}.
\]
On peut encore, avec les notations de
(\hyperref[III.1.2.2-fr]{1.2.2}), identifier
$M_{f_{i_0}f_{i_1}\cdots f_{i_p}}$ à
$\varinjlim_n M_{i_0i_1\ldots i_p}^{(n)}$. Cette identification est un
isomorphisme de $S$-modules gradués, si l'on définit le degré d'un élément
homogène $z\in\varinjlim_n M_{i_0i_1\ldots i_p}^{(n)}$ de la façon
suivante~: $z$ est l'image canonique d'un élément homogène
$x\in M_{i_0i_1\ldots i_p}^{(n)}=M$, de degré $m$~; on prend alors pour
degré de $z$ le nombre
$m-n(d_{i_0}+d_{i_1}+\cdots+d_{i_p})$. Tenant compte de la définition des
homomorphismes
$\varphi_{kn}:M_{i_0i_1\ldots i_p}^{(n)}\to
M_{i_0i_1\ldots i_p}^{(k)}$ (\hyperref[III.1.2.2-fr]{1.2.2}), on voit
aussitôt que cette définition ne dépend pas du «~représentant~» $x$ de $z$
que l'on a considéré. Désignant comme dans
(\hyperref[III.1.2.2-fr]{1.2.2}) par $C_n^p(M)$ l'ensemble des applications
alternées de $[1,r]^{p+1}$ dans $M$ (pour tout $n$), on définit de la même
manière que ci-dessus une structure de $S$-module gradué sur
$\varinjlim_n C_n^p(M)$~; on a encore comme dans
(\hyperref[III.1.2.2-fr]{1.2.2})
\begin{equation}
\label{III.2.1.2.2-fr}
\tag{2.1.2.2}
C^p(\mathfrak U,\mathcal F(*))=\varinjlim_n C_n^p(M)
\end{equation}
l'isomorphisme des deux membres \emph{respectant les degrés}. On a alors,
comme dans (\hyperref[III.1.2.2-fr]{1.2.2})
\begin{equation}
\label{III.2.1.2.3-fr}
\tag{2.1.2.3}
C^p(\mathfrak U,\mathcal F(*))
=C^{p+1}((\mathbf f),M)
=\varinjlim_n K^{p+1}(\mathbf f^n,M)
\end{equation}
l'isomorphisme conservant les degrés~: le degré d'un élément de
$\varinjlim_n K^{p+1}(\mathbf f^n,M)$, image canonique d'une cochaîne
$g\in K^{p+1}(\mathbf f^n,M)$ dont les valeurs
$g(i_0,\ldots,i_p)$ sont dans une même composante homogène $M_m$ de $M$,
est $m-n(d_{i_0}+\cdots+d_{i_p})$, et il est indépendant du choix de cette
cochaîne comme représentant de l'élément considéré.
\end{env}

Comme les isomorphismes précédents sont compatibles avec les opérateurs
cobords, on en conclut, comme dans
(\hyperref[III.1.2.2-fr]{1.2.2}) que l'on a~:

\begin{proposition}[2.1.3]
\label{III.2.1.3-fr}
Soit $X$ un préschéma dont l'espace sous-jacent est noethérien, ou un schéma
quasi-compact. Il existe un isomorphisme canonique fonctoriel en $\mathcal F$
\begin{equation}
\label{III.2.1.3.1-fr}
\tag{2.1.3.1}
H^p(\mathfrak U,\mathcal F(*))
\xrightarrow{\sim}H^{p+1}((\mathbf f),M)
\qquad\text{pour tout }p\geq1.
\end{equation}
On a en outre une suite exacte fonctorielle en $\mathcal F$
\begin{equation}
\label{III.2.1.3.2-fr}
\tag{2.1.3.2}
0\longrightarrow H^0((\mathbf f),M)\longrightarrow M
\longrightarrow H^0(\mathfrak U,\mathcal F(*))
\longrightarrow H^1((\mathbf f),M)\longrightarrow0.
\end{equation}
De plus, tous les homomorphismes introduits sont de degré $0$ pour les
structures de $S$-modules gradués ($S$ étant l'anneau
(\hyperref[III.2.1.1.1-fr]{2.1.1.1})).
\end{proposition}
\providecommand{\bbGamma}{\boldsymbol{\Gamma}}
\providecommand{\shProj}{\mathcal{P}\mathrm{roj}}

